\documentclass{article}

\usepackage{packages}
\usepackage{macros}

\title{Geometric Type Theory}
\author{Johannes Schipp von Branitz \and Mark Damuni Williams}

\begin{document}
\maketitle


\section{Introduction}

\section{Axioms}

\begin{enumerate}
  \item There is a topos $\cB$.
        \textbf{Remark}: It would be better to say there is a logos $\cB$. 
        We should imagine $\cB$ comes with algebraic operations; $\mathrm{colim}$ and $\mathrm{lim}$. with corresponding algebraic laws. Some difficulty with this includes presentability.. Maybe this doesn't matter so much internally? 

  \item \todo{Find a property of $\cB$ corresponding to the \'etale topology}
  \item For every (finitely/countably) presented geometric theory $\bT$, the natural map
    \[ \cB[\bT] \to (\mathrm{Geom}_{\cB}(\cB[\bT], \cB) \to \cB) \]
  is an equivalence.
  \item \'Etale locale choice.
\end{enumerate}

We introduce notation $\mathrm{Sh}(\mathcal{L}) := \mathrm{Log}_{\cB}(\mathcal{L}, \cB)$. This corresponds to our algebraic/geometric intuition. 
$\cB$ is our base logos, and  $\mathrm{Sh}(\cB) = 1$.
$\cB[X]$ is the logos of sets, and $\mathrm{Sh}(\cB[X]) = \cB$.  

\section{Realisation}

Given an object $x : \cB$ there is a realisation $T(x)$ given by $\hom(1, x)$. 
For instance, take the natural numbers object $N : \cB$.
Then we have a map $\bN \to T(N)$ given by, for each $n$ taking the internal successor of zero $n$-times. 

We have a universe $\UU_\cB := \Sigma_{x : \cB} \hom(1, x)$.
\begin{lemma}
  The universe $\UU_\cB$ is closed under $\Pi$, $\Sigma$, $=$, limits/colimits, $\bN$.
\end{lemma}


More generally given a logos $\cB \to \mathcal{L}$ and an $x : \mathcal{L}$ we 
have a realisation 

\section{Simplicial aspects}

There are several options for what could be used as an ``interval''.
There is $\mathcal{S} := [\mathbf{2}, \cB]$ the ``sierpinski'' topos. 
$\cB$-logos morphisms from $Sh(\mathcal{S}) := \mathcal{S} \to \cB$ correspond to objects of $\cB$ with at most $1$ element. 
That is $\mathrm{Sh}(\mathcal{S}) = \Omega \subseteq \cB$ is the geometric propositions. 



\end{document}